% Coordenadas curvilíneas: base q̂_i emanando de P (libro fig. 3.8)
\documentclass[border=5pt]{standalone}
\usepackage{tikz}
\usetikzlibrary{arrows.meta, calc}
\begin{document}
\begin{tikzpicture}[>={Stealth[length=2.4mm]}, line width=0.7pt]
  \coordinate (O) at (0,0);
  \draw[dashed] (O)--(-1.7,-1.05) node[below left]{$x$};
  \draw[dashed] (O)--(2.9,0) node[right]{$y$};
  \draw[dashed] (O)--(0,3.2) node[above]{$z$};
  \coordinate (P) at (1.4,1.55);
  \draw[dashed] (O)--(P);
  \fill (P) circle (1.3pt);
  \node[left=2pt] at (P){$P(q_1,q_2,q_3)$};
  \draw[->, line width=1.1pt] (P)--($(P)+(0.25,0.85)$) node[above]{$\hat q_1$};
  \draw[->, line width=1.1pt] (P)--($(P)+(0.85,0.18)$) node[right]{$\hat q_2$};
  \draw[->, line width=1.1pt] (P)--($(P)+(0.6,-0.6)$) node[below right]{$\hat q_3$};
\end{tikzpicture}
\end{document}
